\begin{theorem}[Scheide Proposition 4.2]
Let $G$ satisfy the paper's criticality convention, here expressed as
edge-criticality together with the requirement that every vertex is
incident with an edge.  Suppose $\chi'(G)=k+1$ and
$\Delta(G)+1\leq k$.  If
\[
 h(G)>t(G)-4,
\]
then $G$ is elementary.
\end{theorem}
\begin{proof}
By \noderef{GlobalMaximumTripleExists}, choose a globally maximum triple.
Then \noderef{ScheideBalancedExists} supplies a balanced normal globally
maximum triple $(T,e_0,\varphi)$; write $h=h(G)$ and $t=t(G)$.

The balanced condition identifies the normal height with $h$.  Its normal
height is at most $\operatorname{length}(T)+1$ by
\noderef{NormalGlobalMaximumTriple}.  On the other hand,
\noderef{GlobalMaximumTriple} and vertex injectivity in
\noderef{TreeSequence} give
$\operatorname{length}(T)+1=|V(T)|=t$.  Hence $h\leq t$.
The node \noderef{ScheideGraphHeightOdd}, applied to the initially chosen
globally maximum triple, says that $h$ is odd.  The maximum-order fact
\noderef{ScheideGlobalMaximumOrder}, applied to the balanced triple, says
that $t$ is odd.

The strict inequality $h>t-4$ says that the nonnegative difference
$t-h$ is less than four.  Since $h$ and $t$ have the same parity, this
difference is even.  Therefore it is either zero or two; equivalently,
\[
 h=t \qquad\text{or}\qquad h+2=t.
\]
In the first case, all the hypotheses of
\noderef{ScheideLargeHeightFullTrunk} hold for the chosen balanced triple,
so that lemma makes $G$ elementary.  In the second case,
\noderef{ScheideLargeHeightOneTailPair} applies to the same triple and
again makes $G$ elementary.  These two exhaustive parity cases prove the
proposition.
\end{proof}
